The study of programming languages is ultimately the study of a class of \textbf{formal systems}.
The programming language itself is just an \textit{implementation} of a formal system.
As such, we need to concern ourselves at least briefly with these notions.

Let's paint a picture.
When we program we type a bunch of symbols into a computer which get saved to a file.
Depending on what language we're using (including what Human language we're assuming) there is a fixed number of characters we're allowed to type into a file.
We call this set of symbols the \textbf{alphabet}.
Of the possible orientations of symbols we can type into a file, only a handful of them are going to be well-formed programs.
Of the well-formed programs, only a handful those make sense as programs.
It doesn't, for example, make sense to add a string to an integer; in other words, only a handful of well-formed programs are well-typed.
And of the well-typed programs, only a handful of \textit{those} make sense as computations.
the program
\begin{ocaml}
  let rec f x = f x
\end{ocaml}
Make sense as an OCaml program in the sense that we can \textit{try} to compute its value, but we'll never succeed (it's an infinite loop).

\todo{A picture, nested squares}

We're interested in formal systems of this form.  In broad strokes, a formal system is made up of a formal language, determined by a \textbf{formal grammar} and one or more \textbf{deductive systems}.
At each level, we have a deductive system which describes a subsection of system.
If each step is \textit{decidable} or \textit{semi-decidable}, meaning we can determine which satisfy what, then we can implement it as a programming language.

\section{Inference Rules}

Deductive systems are formalizations of deductive reasoning over what are called \textbf{judgments}.
They are made up of \textbf{inference rules}, which can be though of as functions which define when we can derive a new judgment from a collection of old ones.

\begin{definition}
  A \textbf{inference rule} is...
\end{definition}

\begin{example}

\end{example}

When we define an inference rule, we're often defining a \textit{collection} of inference rules simultaneously.
This is done by the use of \textbf{meta-variables}.
A meta-variable stands for a syntactic construct whose form is determined by context.

One we have a notion of inference rule, we need the notion of a \textbf{derivation}.

\begin{definition}
A \textbf{derivation} is a tree with the following properties.
\end{definition}

\section{Derivations}
